\documentclass[a4paper, serif, 10pt]{moderncv}

%% moderncv
% style and color
\moderncvstyle{banking}
\moderncvcolor{black}

%% page layout/margins
\usepackage[top=2.5cm, bottom=2.5cm, left=1.5cm, right=1.5cm]{geometry}

\nopagenumbers{}

%% Babel config for Indonesian (Bahasa Indonesia) language
% ref:
% - https://latex3.github.io/babel/guides/locale-indonesian.html
\usepackage[indonesian]{babel}

%% fontspec
\usepackage{fontspec}

\IfFontExistsTF{Linux Libertine}{
  \setmainfont[Ligatures={TeX, Common}, Numbers=OldStyle]{Linux Libertine}
}{}
\IfFontExistsTF{Linux Libertine O}{
  \setmainfont[Ligatures={TeX, Common}, Numbers=OldStyle]{Linux Libertine O}
}{}

%% CTeX
% CJK support, used to show CJK characters in the resume
%
% - fontset=none: disable builtin fontset but instead set the CJK font manually
% - heading=false: disable ctex heading
% - punct=kaiming: use kaiming punctuations styles for CJK
% - scheme=plain: use plain scheme, do not override \normalsize font size
% - space=auto: space settings for CJK characters
%
% ref:
% - http://ctan.mirrorcatalogs.com/language/chinese/ctex/ctex.pdf
\usepackage[UTF8, heading=false, punct=kaiming, scheme=plain, space=auto]{ctex}

\IfFontExistsTF{Noto Serif CJK SC}{
  \setCJKmainfont{Noto Serif CJK SC}
}{}
\IfFontExistsTF{Noto Sans CJK SC}{
  \setCJKsansfont{Noto Sans CJK SC}
}{}

%% line spacing
\usepackage{setspace}
\setstretch{1.125}

%% URL styling - use normal text instead of monospace
\urlstyle{same}

% Auto-underline all links
% Defer until \begin{document} so hyperref is already loaded
\AtBeginDocument{%
  \let\oldhref\href
  \renewcommand{\href}[2]{\oldhref{#1}{\underline{#2}}}%
}

%% Patch \httplink and \httpslink to handle full URLs
% The original moderncv macros always prepend http:// or https:// to URLs.
% This causes issues when the URL already has a protocol (e.g., https://example.com)
% resulting in malformed URLs like https://https://example.com.
% This patch checks if the URL already contains "://" and uses it directly if so.
% Note: We use \str_set:Nx (x-expansion) to expand macros like \@homepage first.
\ExplSyntaxOn
\str_new:N \g_yamlresume_url_str

\RenewDocumentCommand{\httpslink}{O{}m}{
  \str_set:Nx \g_yamlresume_url_str {#2}
  \str_if_in:NnTF \g_yamlresume_url_str {://}
    { \url{#2} }
    { \url{https://#2} }
}

\RenewDocumentCommand{\httplink}{O{}m}{
  \str_set:Nx \g_yamlresume_url_str {#2}
  \str_if_in:NnTF \g_yamlresume_url_str {://}
    { \url{#2} }
    { \url{http://#2} }
}
\ExplSyntaxOff

\name{Sari Wulandari}{}
\title{UX Designer Berfokus pada Riset Pengguna dan Desain Inklusif}
\phone[mobile]{+62 812-3456-7890}
\email{sari.wulandari@email.com}
\homepage{https://sariwulandari.design}

\address{Jl. Jenderal Sudirman No. 12, Jakarta Selatan, DKI Jakarta, Indonesia, 12190}{}{}

\extrainfo{{\small \faLinkedin}\ \href{https://www.linkedin.com/in/sariwulandari}{@sariwulandari} {} {} {} • {} {} {}
{\small \faBehance}\ \href{https://www.behance.net/sariwulandari}{@sariwulandari}}

\begin{document}

\maketitle

\section{Informasi Dasar}

\cvline{}{\begin{itemize}
\item UX Designer dengan lebih dari 6 tahun pengalaman merancang produk digital yang berpusat pada pengguna untuk industri fintech, kesehatan, dan e-commerce.
\item Terampil dalam menjalankan riset pengguna, usability testing, wireframing, prototyping, dan membangun design system yang konsisten.
\item Berpengalaman berkolaborasi dengan product manager, engineer, dan stakeholder untuk menghadirkan solusi yang memenuhi kebutuhan bisnis dan pengguna.
\item Memiliki pemahaman kuat tentang prinsip aksesibilitas dan desain inklusif, serta terbiasa menggunakan Figma, Sketch, dan tools riset seperti Maze dan Hotjar.
\end{itemize}}
\section{Pendidikan}

\cventry{Agu 2012–Jul 2016}
        {Sarjana, Desain Produk, Nilai: 3.78/4.00}
        {Institut Teknologi Bandung}
        {\url{https://www.itb.ac.id}}
        {}
        {\begin{itemize}
\item Mempelajari dasar-dasar desain produk, termasuk ergonomi, estetika, dan pengalaman pengguna.
\item Menyelesaikan proyek akhir tentang aplikasi mobile untuk lansia, yang menjadi fondasi minat pada desain inklusif.
\item Aktif dalam organisasi mahasiswa desain dan menjadi asisten laboratorium untuk mata kuliah interaksi manusia dan komputer.
\end{itemize}
\textbf{Mata Kuliah}: Interaksi Manusia dan Komputer, Psikologi Desain, Desain Antarmuka Digital, Metode Riset Desain, Desain Berbasis Bukti, Manajemen Proyek Desain}
\section{Pengalaman Kerja}

\cventry{Jan 2022–Sekarang}
        {Senior UX Designer}
        {PT Finansial Nusantara}
        {\url{https://www.finansialnusantara.co.id}}
        {}
        {\begin{itemize}
\item Memimpin proses desain end-to-end untuk aplikasi mobile perbankan dengan 5 juta pengguna aktif.
\item Membangun dan mengelola design system bernama "Aksara" yang digunakan oleh 12 tim produk, meningkatkan konsistensi desain hingga 40\%.
\item Melakukan riset pengguna secara berkala, termasuk depth interview dan usability testing dengan responden dari berbagai segmen.
\item Berkolaborasi dengan product manager dan engineer untuk memprioritaskan perbaikan UX berdasarkan data analitik dan feedback pengguna.
\item Mentoring tiga desainer junior dan mengadakan workshop desain internal setiap bulan.
\end{itemize}
\textbf{Kata Kunci}: Riset Pengguna, Sistem Desain, Pengujian Kegunaan, Perbankan Mobile, Mentoring Desain, Figma}

\cventry{Mar 2019–Des 2021}
        {Product Designer}
        {PT Kesehatan Digital Indonesia}
        {\url{https://www.kesehatandigital.id}}
        {}
        {\begin{itemize}
\item Merancang pengalaman pengguna untuk aplikasi telemedicine dan platform manajemen rekam medis.
\item Menyusun user journey map dan service blueprint untuk mengidentifikasi titik perbaikan layanan.
\item Mengembangkan prototipe interaktif menggunakan Figma dan menjalankan usability testing dengan 100+ partisipan.
\item Meningkatkan tingkat penyelesaian alur pendaftaran pasien dari 58\% menjadi 82\% melalui penyederhanaan form dan microcopy.
\item Menjalin komunikasi erat dengan tim klinis dan engineering untuk memastikan kepatuhan terhadap regulasi kesehatan.
\end{itemize}
\textbf{Kata Kunci}: Telemedicine, Peta Perjalanan Pengguna, Pembuatan Prototipe, Tulisan Antarmuka, Regulasi Kesehatan}

\cventry{Sep 2016–Feb 2019}
        {Junior UX Designer}
        {PT Media Kreasi Nusantara}
        {\url{https://www.mediakreasi.co.id}}
        {}
        {\begin{itemize}
\item Mendukung desain antarmuka untuk portal berita dan aplikasi media sosial internal.
\item Membantu pelaksanaan user testing dan menyusun laporan temuan untuk tim produk.
\item Membuat wireframe dan mockup untuk fitur-fitur baru dengan pendekatan mobile-first.
\item Berpartisipasi dalam sprint planning dan daily stand-up untuk memastikan kelancaran iterasi desain.
\end{itemize}
\textbf{Kata Kunci}: Wireframe, Pengujian Pengguna, Mobile-First, Agile, Sketch}
\section{Bahasa}

\cvline{Indonesia}{Bahasa Ibu / Bilingual \hfill \textbf{Kata Kunci}: Bahasa sehari-hari dan formal}
\cvline{Inggris}{Tingkat Profesional Penuh \hfill \textbf{Kata Kunci}: TOEFL ITP 620, IELTS 7.0}
\section{Keahlian}

\cvline{Riset Pengguna}{Ahli \hfill \textbf{Kata Kunci}: Wawancara Mendalam, Survei, Analisis Kompetitor, Persona, Peta Perjalanan Pengguna}
\cvline{Desain Interaksi}{Ahli \hfill \textbf{Kata Kunci}: Wireframe, Pembuatan Prototipe, Desain Interaksi, Interaksi Mikro, Aksesibilitas}
\cvline{Sistem Desain}{Mahir \hfill \textbf{Kata Kunci}: Token Desain, Komponen UI, Dokumentasi, Pustaka Figma}
\cvline{Pengujian dan Analitik}{Mahir \hfill \textbf{Kata Kunci}: Pengujian Kegunaan, Pengujian A/B, Hotjar, Maze, Google Analytics}
\section{Penghargaan}

\cventry{Nov 2023}
        {Asosiasi UX Indonesia}
        {Penghargaan Desain Inklusif 2023}
        {}
        {}
        {Diberikan kepada perancang yang berhasil menghadirkan solusi digital yang dapat diakses oleh pengguna dengan berbagai kemampuan, termasuk lansia dan penyandang disabilitas.}
\section{Sertifikat}

\cventry{Agu 2021}
        {Google}
        {Google UX Design Professional Certificate}
        {\url{https://www.coursera.org/professional-certificates/google-ux-design}}
        {}
        {}
\section{Publikasi}

\cventry{Jun 2023}
        {Meningkatkan Kepercayaan Pengguna melalui Desain Transparan pada Aplikasi Finansial}
        {Medium}
        {\url{https://medium.com/@sariwulandari/meningkatkan-kepercayaan-pengguna}}
        {}
        {Artikel ini membahas pentingnya transparansi informasi dan kontrol pengguna dalam aplikasi finansial, dilengkapi studi kasus dan rekomendasi praktis untuk desainer produk.}
\section{Proyek}

\cventry{Apr 2020–Nov 2021}
        {Aplikasi mobile untuk konsultasi dokter dan pemantauan kesehatan pasien.}
        {SehatKu}
        {\url{https://www.sehatku.id}}
        {}
        {\begin{itemize}
\item Merancang alur konsultasi online yang intuitif, termasuk pemilihan dokter, jadwal, dan pembayaran dalam satu aplikasi.
\item Membangun komponen desain yang dapat digunakan ulang untuk mempercepat pengembangan fitur baru.
\item Berkolaborasi dengan tim medis untuk memastikan informasi kesehatan disajikan secara jelas dan tidak menyesatkan.
\end{itemize}
\textbf{Kata Kunci}: Kesehatan, Aplikasi Mobile, Telemedicine, Desain Inklusif}
\section{Minat}

\cvline{Membaca}{Buku desain, Fiksi sains, Jurnal UX}
\cvline{Fotografi}{Street photography, Komposisi visual, Pengeditan foto}
\cvline{Komunitas}{UX Indonesia, Mentoring, Workshop desain}
\section{Kegiatan Sukarela}

\cventry{Jan 2022–Des 2023}
        {Mentor Desain}
        {Komunitas UX Indonesia}
        {\url{https://www.uxindo.id}}
        {}
        {\begin{itemize}
\item Menjadi mentor untuk desainer muda dan profesional yang ingin beralih ke bidang UX.
\item Menyusun materi lokakarya tentang riset pengguna, design thinking, dan penyusunan portofolio.
\item Membantu menyelenggarakan acara tahunan "UX Fair" yang dihadiri lebih dari 500 peserta.
\end{itemize}}

\end{document}